\documentclass{article}

\usepackage{enumitem}
\usepackage{listings}
\usepackage{color}
\usepackage{amsmath}
\usepackage{hyperref}
\usepackage{graphicx}
\usepackage{pgffor}
\usepackage{xparse}
\usepackage{expl3}
\usepackage{tabularx, makecell}
\usepackage{multirow}
\usepackage{booktabs}
\usepackage{indentfirst}
\usepackage{lipsum}
\usepackage{sectsty}
\usepackage[utf8]{inputenc}
\usepackage{csquotes}
\usepackage{xcolor}
\usepackage{fancyvrb}
\usepackage{fancyhdr}
\usepackage{fancyvrb}
\usepackage[most]{tcolorbox}
\usepackage{blindtext}
\usepackage{caption}
\usepackage{etoolbox}
\usepackage{booktabs}
\usepackage{karnaugh-map}
\usepackage{tikz}
\usepackage{mdframed}
\usepackage{calc}
\usepackage[nomessages]{fp}
\usepackage{pgfplots}
\usepackage{geometry}
\usepackage{python}
\usepackage{float}
\usepackage{colortbl}

\graphicspath{{./}}

\definecolor{codegreen}{rgb}{0,0.6,0}
\definecolor{codegray}{rgb}{0.5,0.5,0.5}
\definecolor{codepurple}{rgb}{0.58,0,0.82}
\definecolor{backcolour}{rgb}{0.95,0.95,0.92}

\sectionfont{\bfseries\Large\center} 

\geometry{
	a4paper,
	total={170mm,257mm},
	left=20mm,
	top=20mm,
}

\pgfplotsset{compat=1.13}

\begin{document}
	% Vars
	\newcommand{\ROne}{99.8}
	\newcommand{\RTwo}{198}
	\newcommand{\RThree}{53.7}
	\newcommand{\RMax}{473}
	
	\newcommand{\IPosOne}{42.1}
	\newcommand{\IPosTwo}{96.0}

	\newcommand{\UExOneOne}{4.24}
	\newcommand{\UExOneTwo}{8.43}
	\newcommand{\UExOneThree}{2.26}

	\newcommand{\IExTwoOne}{61.7}
	\newcommand{\IExTwoTwo}{43.2}
	\newcommand{\IExTwoThree}{18.6}

	\newcommand{\UExTwoOne}{6.66}
	\newcommand{\UExTwoTwo}{8.69}
	\newcommand{\UExTwoThree}{8.69}

	% <Square cases>
	\makeatletter
	\newenvironment{sqcases} {
		\matrix@check\sqcases\env@sqcases
	}{
		\endarray \right.
	}
	\def\env@sqcases {
		\let \@ifnextchar \new@ifnextchar
		\left \lbrack
		\def \arraystretch{1.2}
		\array{@{}l@{\quad}l@{}}
	}
	\makeatother
	% </Square cases>

	% My include img. (With frame and caption) 
	\newcommand{\myIncludeImg}[4]{
		\begin{center} \begin{figure}[H]
			\begin{mdframed}[backgroundcolor=black!5, rightline=false, leftline=false]
				\begin{center}
					\includegraphics[scale=#1]{#2}
					\caption{#3}
				\end{center}
			\end{mdframed}
			\label{fig:#4}
		\end{figure} \end{center}
	}

	% Begin of the document

	\title{CDE laboratory\_01}
	\author{Terman Emil FAF161}
	\maketitle

	% Write at bottom
	\vspace*{\fill}
	
	\centering
	\includegraphics{imgs/UTM_logo.png}

	\begin{flushright}
		Prof: O. Lupan
	\end{flushright}

	\LaTeX
	\pagebreak %-----------------------------------------------End of title page

	\raggedright
	\section{Verification of Kirchhoff and Ohm's laws for electrical circuits.}
		\textbf{Kirchhoff first law:} The current entering any junction is equal to the current leaving that junction.
		\textbf{Kirchhoff second law:} The sum of all the voltages around a loop is equal to zero.

		\myIncludeImg{0.7}{./imgs/ElectricCircuit1.jpg}{Circuit 1.1}{circuit1}

		\subsection{}
			\[
				\begin{tabular}{|c|c|c|c|c|c|c|c|}
					\hline
					$R_1 (\Omega)$ & $R_2 (\Omega)$ & $R_3 (\Omega)$ & $I_1 (mA)$ & $I_2 (mA)$ & $U_{t1} (V)$ & $U_{t2} (V)$\\
					\hline
					\ROne & \RTwo & \RThree & \IPosOne & \IPosTwo& 15.01 & 15.01\\
					\hline
				\end{tabular}
			\]

		\subsection{}
			Internal resistance:
			\[
				r_0 = \frac{U_{t2} - U_{t1}}{I_1 - I_2} = \frac{0}{I_1 - I_2} = 0
			\]
			The tools we had, weren't exact enough, that's why we got $r_0$ = 0. To find a more exact value of internal resistance, we need something better.

		\subsection{}
			The current in circuit and the tensions in each resistance by the Ohm's Laws.
			% Calculates the total sum of resistances
			\FPeval{\resistanceSum}{clip(\ROne + \RTwo + \RThree)}

			% 15 / (sum of resistance)
			\FPeval{\currentExThree}{clip(15 / \resistanceSum * 1000)}

			% Round \currentExThree to 2 fractional digits
			\FPround{\currentExThree}{\currentExThree}{2}

			\[
				I = \frac{E}{R_1 + R_2 + R_3 + r_0} = \frac{15}{\resistanceSum} = \currentExThree \ mA
			\]
			
			\FPeval{\UOneCalc}{clip(\currentExThree * \ROne / 1000)}
			\FPround{\UOneCalc}{\UOneCalc}{2}

			\FPeval{\UTwoCalc}{clip(\currentExThree * \RTwo / 1000)}
			\FPround{\UTwoCalc}{\UTwoCalc}{2}

			\FPeval{\UThreeCalc}{clip(\currentExThree * \RThree / 1000)}
			\FPround{\UThreeCalc}{\UThreeCalc}{2}

			\[U_1 = IR_1 = \currentExThree \cdot \ROne = \UOneCalc \ V\]
			\[U_2 = IR_2 = \currentExThree \cdot \RTwo = \UTwoCalc \ V\]
			\[U_3 = IR_3 = \currentExThree \cdot \RThree = \UThreeCalc \ V\]

		\subsection{}
			\begin{center} \begin{tabular}{|c|c|c|c|c|c|c|c|c|c|}
				\hline
				\multicolumn{2}{|c|}{R ($\Omega$)} &$
					I_{c}$(mA) &
					\multicolumn{2}{c|}{$U_c$ (V)} &
					$I_{m}$ (mA) &
					\multicolumn{2}{c|}{$U_m$ (V)} \\

				\hline
				
				$R_1$ & \ROne &
					\multirow{3}{*}{\currentExThree} &
					$U_1$ & \UOneCalc &
					\multirow{3}{*}{\IPosOne} &
					$U_1$ & \UExOneOne\\
				
				\cline{1-2} \cline{4-5} \cline{7-8}

				$R_2$ & \RTwo & & $U_2$ & \UTwoCalc & & $U_2$ & \UExOneTwo\\		
				\cline{1-2} \cline{4-5} \cline{7-8}
				
				$R_3$ & \RThree & & $U_3$ & \UThreeCalc & & $U_3$ & \UExOneThree\\
				\hline
			\end{tabular} \end{center}

		\subsection{Verify the second Kirchhoff's law}
			\FPeval{\UTotalCalc}{clip(\UOneCalc + \UTwoCalc + \UThreeCalc)}
			\FPeval{\UTotalMes}{clip(\UExOneOne + \UExOneTwo + \UExOneThree)}

			\[U_{c1} + U_{c2} + U_{c3} = \UOneCalc + \UTwoCalc + \UThreeCalc = \textbf{\UTotalCalc \ V}\]
			\[U_{m1} + U_{m2} + U_{m3} = \UExOneOne + \UExOneTwo + \UExOneThree = \textbf{\UTotalMes \ V}\]

		\subsection{}
			\myIncludeImg{0.3}{./imgs/PotentialChart.jpg}{Potential Chart of Circuit 1}{potentialChart}

		\subsection{Verify the first law of Kirchhoff}
			\myIncludeImg{0.7}{./imgs/ElectricCircuit2.jpg}{Circuit 1.2}{circuit2}
			
			The equivalent resistance of the circuit $R_E$:
			\FPeval{\REtotal}{clip(\ROne + \RTwo * \RMax / (\RTwo + \RMax))}
			\FPround{\REtotal}{\REtotal}{2}
			\[
				R_E = R_1 + \frac{R_2 R_3}{R_2 + R_3} = \ROne + \frac{\RTwo \cdot \RMax}{\RTwo + \RMax} = \REtotal \ \Omega
			\]

			\FPeval{\IExTwoCalcOne}{clip(15 / \REtotal * 1000)}
			\FPround{\IExTwoCalcOne}{\IExTwoCalcOne}{2}

			\[
				I_1 = \frac{E}{r_0 + R_E} = \frac{15}{0 + \REtotal} = \IExTwoCalcOne \ mA
			\]

			\FPeval{\UExTwoCalcTwo}{clip(\IExTwoCalcOne * \RTwo * \RMax / (\RTwo + \RMax) / 1000)}
			\FPround{\UExTwoCalcTwo}{\UExTwoCalcTwo}{2}

			\[
				U_2 = U_3 = I_1 \frac{R_2 R_3}{R_2 + R_3} =
					\IExTwoCalcOne \cdot \frac{\RTwo \cdot \RMax}{\RTwo + \RMax} =
					\UExTwoCalcTwo \ V
			\]

			\FPeval{\UExTwoCalcOne}{clip(\IExTwoOne * \ROne / 1000)}
			\FPround{\UExTwoCalcOne}{\UExTwoCalcOne}{2}

			\[
				U_1 = I_1 \cdot R_1 = \IExTwoOne \cdot \ROne = \UExTwoCalcOne \ V
			\]

			%I_1
			\FPeval{\IExTwoCalcTwo}{clip(\UExTwoCalcTwo / \RTwo * 1000)}
			\FPround{\IExTwoCalcTwo}{\IExTwoCalcTwo}{2}

			%I_2
			\FPeval{\IExTwoCalcThree}{clip(\UExTwoCalcTwo / \RMax * 1000)}
			\FPround{\IExTwoCalcThree}{\IExTwoCalcThree}{2}

			\begin{center} \begin{tabular}{c|c}
				\(I_2 = \frac{U_2}{R_2} = \frac{\UExTwoCalcTwo}{\RTwo} = \IExTwoCalcTwo \ mA\)
				&
				\(I_3 = \frac{U_3}{R_3} = \frac{\UExTwoCalcTwo}{\RMax} = \IExTwoCalcThree \ mA\)\\
			\end{tabular} \end{center}

			\begin{center} \begin{tabular}{|c|c|c|c|c|c|c|c|c|c|}
				\hline
					\multicolumn{2}{|c|}{R ($\Omega$)} &
					\multicolumn{2}{c|}{$I_c (mA)$} &
					\multicolumn{2}{c|}{$U_c (V)$} &
					\multicolumn{2}{c|}{$I_m (mA)$} &
					\multicolumn{2}{c|}{$U_m (V)$}\\
				\hline

				$R_1$ & \ROne & $I_1$ & \IExTwoCalcOne & $U_1$ & \UExTwoCalcOne & $I_1$ & \IExTwoOne & $U_1$ & \UExTwoOne\\
				\hline

				$R_2$ & \RTwo & $I_2$ & \IExTwoCalcTwo & $U_2$ & \UExTwoCalcTwo & $I_2$ & \IExTwoTwo & $U_2$ & \UExTwoTwo\\
				\hline

				$R_3$ & \RMax & $I_3$ & \IExTwoCalcThree & $U_3$ & \UExTwoCalcTwo & $I_3$ & \IExTwoThree & $U_3$ & \UExTwoThree\\
				\hline
			\end{tabular} \end{center}
		
		\subsection{}
			\FPeval{\IExTwoKirchhoffDiff}{clip(\IExTwoTwo + \IExTwoThree - \IExTwoOne)}
			\FPround{\IExTwoKirchhoffDiff}{\IExTwoKirchhoffDiff}{2}

			\begin{minipage}{\textwidth}
				\begin{center} \textbf{Kirchhoff's first Law} \end{center}
				\[
					\sum_{1}^{n} I_k = 0 \Rightarrow I_2 + I_3 - I_1 = 0 \Rightarrow \IExTwoTwo + \IExTwoThree - \IExTwoOne = \IExTwoKirchhoffDiff \approx 0
				\]
				\hrule
				\begin{center} \textbf{Power Balance} \end{center}

				%Left part
				\FPeval{\ExTwoPowerEqL}{clip(15 * \IExTwoOne)}
				\FPround{\ExTwoPowerEqL}{\ExTwoPowerEqL}{2}

				%Right part
				\FPeval{\ExTwoPowerEqR}{clip(\IExTwoOne * \IExTwoOne * (0 + \ROne))}
				\FPeval{\ExTwoPowerEqR}{clip(\ExTwoPowerEqR + \IExTwoTwo * \IExTwoTwo * \RTwo)}
				\FPeval{\ExTwoPowerEqR}{clip(\ExTwoPowerEqR + \IExTwoThree * \IExTwoThree * \RMax)}
				\FPeval{\ExTwoPowerEqR}{clip(\ExTwoPowerEqR / 1000)}
				\FPround{\ExTwoPowerEqR}{\ExTwoPowerEqR}{2}

				\[
					E \cdot I_1 = I_1 ^ 2 (r_0 + R_1) + I_2 ^ 2 R_2 + I_3 ^ 2 R_3 \Rightarrow
				\]
				\[
					15 \cdot \IExTwoOne = \frac{1}{10 ^ 3} \cdot (\IExTwoOne ^ 2 \cdot (0 + \ROne) + \IExTwoTwo ^ 2 \cdot \RTwo + \IExTwoThree ^ 2 \cdot \RMax)  \Leftrightarrow
				\]
				\[
					\ExTwoPowerEqL \ W \approx \ExTwoPowerEqR \ W
				\]
			\end{minipage}


		\subsection{}
			\begin{center} \begin{tabular}{|c|c|c|c|c|c|c|c|c|c|}
				\hline
				\multicolumn{7}{|c}{\textbf{Measured}} & \multicolumn{3}{|c|}{\textbf{Calculated}}\\
				\hline
				$R_3$ & U & $U_1$ & $U_2$ & $I_1$ & $I_2$ & $I_3$ & $U_1 + U_2$ & $I_2 + I_3$ & P\\
				\hline
				$\Omega$ & \multicolumn{3}{c|}{V} & \multicolumn{3}{c|}{mA} & V & mA & mW\\	
				\hline
				0 & \multirow{6}{*}{15} & 14.93 & 0.06 & 144.3 & 0.2 & 143.5 & 14.99 & 143.7 & 2164.5\\
				100 & & 9.06 & 5.88 & 88.4 & 29.4 & 59.2 & 14.94 & 88.6 & 1326\\
				200 & & 7.57 & 7.42 & 74 & 37.1 & 36.7 & 14.99 & 73.8 & 1110\\
				300 & & 6.9 & 8.0 & 67.2 & 40.5 & 26.7 & 14.9 & 67.2 & 1008\\
				467 & & 6.35 & 8.64 & 62 & 43.3 & 18.6 & 14.99 & 61.9 & 930\\
				\hline
			\end{tabular} \end{center}
			\[
				P_{R_3} = U \cdot I_{R_3} 
			\]

			\myIncludeImg{0.34}{./imgs/U1U2_I1I2I3_R3_plot.jpg}{}{graph1}

			\subsection{Conclusion}
				In this laboratory work we showed that the Kirchhoff's laws are true, by making our own circuits. The errors appeared mostly because of the human factor and because of the relatively bad equipment.
				\par In the first graph is shown the second Kirchhoff's law: \(U_1 + U_2 = 15V\)
				\par From the second graph, we can clearly see the first Kirchhoff's Law: \(I_1 = I_2 + I_3\).
				

	\newpage
	\section{De cercetat proprietatile elementelor pasive (R, L, C) în circuitul de curent alternativ.}
		\myIncludeImg{0.7}{./imgs/ElectricCircuit3.jpg}{Circuit 2.1}{circuit3}
		\myIncludeImg{0.7}{./imgs/ElectricCircuit4.jpg}{Circuit 2.2}{circuit4}

		\begin{center} \begin{tabular}{|c|c|c|c|c|c|c|c|c|c|c|c|c|c|}
			\hline
			
			\multirow{2}{*}{Element} & U & $U_m$ & I & $I_m$ & $\phi$ & Q & S & P & R & C & L & $X_C$ & $X_L$\\
			\arrayrulecolor{gray!70}
			\cline{2-14}
			\arrayrulecolor{black}
			
			& \multicolumn{2}{c|}{V} & \multicolumn{2}{c|}{mA} & o & VAR & VA & W & $\Omega$ & nF & mH & $\Omega$ & $\Omega$\\
			\hline
			
			$R = 510 \Omega$ & 10.3 & & 20 &  & & & & & & & & &\\
			\hline
			
			C = 56 nF & 4.43 & & 186.3 & & & & & & & & & &\\
			\hline
			
			L = 3.64 mH & 11.34 & & 510 & & & & & & & & & &\\
			\hline
		\end{tabular} \end{center}

		\myIncludeImg{0.1}{./imgs/Part2OscilographPos1.jpg}{Oscilograph pos 1}{oscil. 1}
		\myIncludeImg{0.1}{./imgs/Part2OscilographPos2.jpg}{Oscilograph pos 2}{oscil. 2}
		\myIncludeImg{0.1}{./imgs/Part2OscilographPos3.jpg}{Oscilograph pos 3}{oscil. 3}
\end{document}
